\documentclass[assignment.tex]{subfiles}
\begin{document}

\chapter{Assignment 1: Syntactic Analysis Pipeline}
\label{sec:assignment1}

Legal contracts contain long, complex sentences with nested clauses, coordinating
structures, and specialised vocabulary.  The first stage of the pipeline decomposes
raw contract text into atomic, analysable units and extracts their syntactic
structure.  Three components are described below: clause splitting, noun-phrase (NP)
chunking, and full dependency parsing.

\section{Clause Splitting}
\label{subsec:clause-splitting}

\subsection{Motivation}

A single legal sentence can span several lines and encode multiple independent
obligations or conditions.  Downstream semantic models (NER, SRL, intent
classification) perform significantly better when each input is a single, coherent
proposition.  The clause splitter reduces complex sentences to their atomic units
while preserving the integrity of conditional constructions.

\subsection{Algorithm}

The splitter operates in five sequential stages:

\begin{enumerate}[leftmargin=*, itemsep=4pt]
    \item \textbf{Paragraph segmentation.}  The raw input text is partitioned on
          double newlines, treating each paragraph as an independent processing unit
          to avoid cross-paragraph sentence boundary errors.

    \item \textbf{Sentence segmentation.}  Each paragraph is processed with
          spaCy's \path{en_core_web_sm} model, which detects sentence boundaries
          more reliably than purely rule-based tokenisers.

    \item \textbf{Conditional detection.}  For each sentence, the algorithm checks
          whether any token simultaneously satisfies \emph{all three} of the
          following: (a) its part-of-speech tag is \texttt{SCONJ}, (b) its
          dependency relation is \texttt{mark}, and (c) its lower-case form appears
          in the conditional keyword set (Table~\ref{tab:split-keywords}).  When
          this condition is met, the entire sentence is retained as a single
          indivisible clause.  Splitting a conditional sentence would sever the
          logical dependency between the protasis and apodosis.

    \item \textbf{Coordinating-conjunction split.}  For sentences that do not
          trigger the conditional guard, the algorithm scans for \texttt{CC} tokens
          whose lower-case form belongs to the coordinate keyword set
          (Table~\ref{tab:split-keywords}).  A split is \emph{only} performed when
          \textbf{both} resulting fragments independently contain a subject--verb
          pair, verified by requiring at least one \texttt{nsubj} or
          \texttt{nsubjpass} dependent and one \texttt{ROOT} or main-verb token on
          each side.

    \item \textbf{Fragment pruning.}  Any resulting clause with two tokens or fewer
          is discarded as a non-informative syntactic fragment.
\end{enumerate}

\begin{table}[H]
    \centering
    \caption{Keyword sets governing clause-splitting behaviour.}
    \label{tab:split-keywords}
    \begin{tabular}{@{}lp{9cm}@{}}
        \toprule
        \textbf{Set} & \textbf{Keywords} \\
        \midrule
        Conditional (\texttt{SCONJ}) &
            \textit{if, unless, provided, whereas, subject, upon, when, whenever,
            although, though, because, since, after, before, until} \\[4pt]
        Coordinate (\texttt{CC}) &
            \textit{and, but, or, nor, yet, so} \\
        \bottomrule
    \end{tabular}
\end{table}

\subsection{Input/Output Examples}

\paragraph{Example 1 — Coordinating split.}

\noindent\textbf{Input:}
\begin{quote}
    \textit{``The Employer shall pay the Employee a monthly salary of \$5,000 and
    the Employee shall complete all assigned tasks by the deadline.''}
\end{quote}

\noindent\textbf{Output} — two independent clauses, one per line:
\begin{enumerate}
    \item \textit{The Employer shall pay the Employee a monthly salary of \$5,000}
    \item \textit{the Employee shall complete all assigned tasks by the deadline}
\end{enumerate}

The conjunction \textit{and} carries the \texttt{CC} tag.  Both resulting fragments
contain independent subject--verb pairs (\textit{Employer/shall pay} and
\textit{Employee/shall complete}), so a split is performed.

\paragraph{Example 2 — Conditional retained as single clause.}

\noindent\textbf{Input:}
\begin{quote}
    \textit{``If the Employee fails to report to work within 3 days without notice,
    the Employer may terminate this Agreement.''}
\end{quote}

\noindent\textbf{Output} — single clause (no split):
\begin{enumerate}
    \item \textit{If the Employee fails to report to work within 3 days without
    notice, the Employer may terminate this Agreement.}
\end{enumerate}

The token \textit{If} carries \texttt{SCONJ} POS and \texttt{mark} dependency
relation, and appears in the conditional keyword set; the sentence is therefore
preserved as an atomic unit.

\section{NP Chunking}
\label{subsec:np-chunking}

\subsection{Method}

Noun-phrase chunking uses spaCy's built-in \texttt{doc.noun\_chunks} iterator,
which returns contiguous base NP spans derived from the dependency parse tree.
These spans are converted to the IOB (\textbf{I}nside--\textbf{O}utside--%
\textbf{B}eginning) tagging scheme defined as follows:

\begin{itemize}[leftmargin=*, itemsep=2pt]
    \item \textbf{B-NP}: the first (or only) token of a noun phrase span.
    \item \textbf{I-NP}: any subsequent continuation token within the same span.
    \item \textbf{O}:    a token that does not belong to any noun phrase.
\end{itemize}

Output is written to \texttt{output/chunks.txt} as tab-separated
\textit{token}$\backslash$t\textit{IOB-tag} pairs, one pair per line, with a blank
line separating successive clauses.

\subsection{IOB Tagging Example}

Table~\ref{tab:iob-example} shows the IOB annotation produced for the clause
\textit{``The contracting parties shall maintain confidentiality.''}.

\begin{table}[H]
    \centering
    \caption{IOB tags for \textit{``The contracting parties shall maintain
    confidentiality.''}}
    \label{tab:iob-example}
    \begin{tabular}{@{}ll@{}}
        \toprule
        \textbf{Token}   & \textbf{IOB Tag} \\
        \midrule
        The              & B-NP \\
        contracting      & I-NP \\
        parties          & I-NP \\
        shall            & O    \\
        maintain         & O    \\
        confidentiality  & B-NP \\
        .                & O    \\
        \bottomrule
    \end{tabular}
\end{table}

The determiner \textit{The}, the pre-modifying participle \textit{contracting}, and
the head noun \textit{parties} together form the subject NP.  The modal auxiliary
\textit{shall} and main verb \textit{maintain} are outside any NP and are tagged
\texttt{O}.  The object NP \textit{confidentiality} begins a new span tagged
\texttt{B-NP}.

\section{Dependency Analysis}
\label{subsec:dependency-parsing}

\subsection{Parser}

Full dependency parsing is performed with spaCy's \texttt{en\_core\_web\_sm} model,
which implements a transition-based arc-eager dependency parser trained on the
English OntoNotes 5 corpus and annotated with Universal Dependencies (UD) relation
labels.

\subsection{Output Format}

Each clause is serialised to a JSON object containing the following top-level
fields:

\begin{itemize}[leftmargin=*, itemsep=2pt]
    \item \texttt{clause\_id} — zero-based integer index of the clause in the
          document.
    \item \texttt{text} — the original clause string before tokenisation.
    \item \texttt{tokens} — an array of token objects, each carrying:
    \begin{itemize}
        \item \texttt{id}: position index within the clause (zero-based).
        \item \texttt{token}: surface form as it appears in the text.
        \item \texttt{lemma}: morphologically normalised form.
        \item \texttt{pos}: universal POS tag (e.g.\ \texttt{NOUN}, \texttt{VERB}).
        \item \texttt{tag}: fine-grained Penn Treebank tag
              (e.g.\ \texttt{NN}, \texttt{VB}, \texttt{MD}).
        \item \texttt{dep}: Universal Dependencies relation label
              (e.g.\ \texttt{nsubj}, \texttt{dobj}, \texttt{ROOT}).
        \item \texttt{head\_id}: position index of the syntactic head token.
        \item \texttt{head\_token}: surface form of the head.
    \end{itemize}
\end{itemize}

The complete output is written to \texttt{output/dependencies.json} as a JSON array
of these clause objects.

\subsection{JSON Output Example}

\begin{lstlisting}[style=jsonStyle,
                   caption={Dependency parse output for the clause
                   \textit{``The Employer shall pay the salary.''}},
                   label={lst:dep-json}]
{
  "clause_id": 0,
  "text": "The Employer shall pay the salary.",
  "tokens": [
    {"id": 0, "token": "The",      "lemma": "the",
     "pos": "DET",   "tag": "DT", "dep": "det",
     "head_id": 1,   "head_token": "Employer"},
    {"id": 1, "token": "Employer", "lemma": "employer",
     "pos": "NOUN",  "tag": "NN", "dep": "nsubj",
     "head_id": 3,   "head_token": "pay"},
    {"id": 2, "token": "shall",    "lemma": "shall",
     "pos": "AUX",   "tag": "MD", "dep": "aux",
     "head_id": 3,   "head_token": "pay"},
    {"id": 3, "token": "pay",      "lemma": "pay",
     "pos": "VERB",  "tag": "VB", "dep": "ROOT",
     "head_id": 3,   "head_token": "pay"},
    {"id": 4, "token": "the",      "lemma": "the",
     "pos": "DET",   "tag": "DT", "dep": "det",
     "head_id": 5,   "head_token": "salary"},
    {"id": 5, "token": "salary",   "lemma": "salary",
     "pos": "NOUN",  "tag": "NN", "dep": "dobj",
     "head_id": 3,   "head_token": "pay"},
    {"id": 6, "token": ".",        "lemma": ".",
     "pos": "PUNCT", "tag": ".",  "dep": "punct",
     "head_id": 3,   "head_token": "pay"}
  ]
}
\end{lstlisting}

The ROOT of the clause is \textit{pay} (index 3).  \textit{Employer} is its nominal
subject (\texttt{nsubj}), \textit{shall} is the auxiliary verb (\texttt{aux}), and
\textit{salary} is the direct object (\texttt{dobj}).  The determiner
\textit{the} (index 0) modifies \textit{Employer} via the \texttt{det} relation,
while \textit{the} (index 4) modifies \textit{salary} similarly.

\end{document}
